\documentclass[11pt,a4paper]{article}

% --- Packages ---
\usepackage[margin=.75in]{geometry}
\usepackage[utf8]{inputenc}
\usepackage[T1]{fontenc}
\usepackage{lmodern} % Fix for font not found error
\usepackage{titlesec}
\usepackage{enumitem}
\usepackage{hyperref}
\usepackage{xcolor}
\usepackage{fancyhdr}
\usepackage{comment}
\usepackage{cleveref}
\usepackage[bottom]{footmisc}
\usepackage{graphicx}
\usepackage{longtable}
\usepackage{booktabs}
\usepackage{array}
\usepackage{calc}
\usepackage{etoolbox}

% Standard Pandoc setup
\providecommand{\tightlist}{%
  \setlength{\itemsep}{0pt}\setlength{\parskip}{0pt}}

\crefformat{section}{\S#2#1#3}
\crefformat{subsection}{\S#2#1#3}
\crefformat{subsubsection}{\S#2#1#3}
\crefformat{figure}{#2Figure~#1#3}
\crefformat{footnote}{#2\footnotemark[#1]#3}

\linespread{0.9} % Riduce lo spazio tra le linee

% --- Colors ---
\definecolor{primary}{RGB}{0, 0, 0}
\definecolor{links}{HTML}{0000EE}
\definecolor{blue}{HTML}{26428B}

\hypersetup{
    colorlinks=true,
    linkcolor=links,
    urlcolor=links,
    citecolor=links,
}

% --- Section Formatting ---
\titleformat{\section}
  {\normalfont\Large\bfseries\color{blue}}{\thesection}{1em}{}
\titleformat{\subsection}
{\normalfont\large\bfseries\color{blue}}{\thesubsection}{1em}{}
\titleformat{\subsubsection}
{\normalfont\bfseries\color{blue}}{\thesubsubsection}{1em}{}

% --- Header and Footer ---
\pagestyle{fancy}
\fancyhead[L]{\small Giuseppe Capasso}
\fancyhead[R]{\small intro-to-linux}
\fancyfoot[C]{\thepage}
\renewcommand{\headrulewidth}{0.4pt}

% --- Custom Title Command ---
\newcommand{\proposaltitle}[1]{
    \begin{center}
        \vspace*{0.1cm}
        {\LARGE \bfseries \color{blue} #1} \\[1.5ex]
        {\large \textbf{Giuseppe Capasso}} \\[1ex]
                \small{
            \vspace{0.1cm}
            \textbf{Materia:} intro-to-linux\\
        }
            \end{center}
    \vspace{0.3cm}
    \hrule height 0.5pt
    \vspace{0.5cm}
}

\begin{document}

\proposaltitle{1. Il File System (FHS): La Mappa del Tesoro (30 min)}

\textbf{Corso:} Linux Livello ITS\\
\textbf{Obiettivo:} Comprendere l'ecosistema Linux, prendere confidenza
con la CLI via browser e preparare l'ambiente di virtualizzazione.

\begin{center}\rule{0.5\linewidth}{0.5pt}\end{center}

\section{1. Il File System (FHS): La Mappa del Tesoro (30
min)}\label{il-file-system-fhs-la-mappa-del-tesoro-30-min}

\textbf{Obiettivo:} Spiegare che non esiste il disco \texttt{C:} e che
\emph{tutto è un file}.

\subsection{Concetti Chiave da
Trasmettere}\label{concetti-chiave-da-trasmettere}

\begin{itemize}
\item
  \textbf{La Radice (\texttt{/})}\\
  L'inizio di tutto l'albero del filesystem.
\item
  \textbf{Mount Point}\\
  Se inserisci una USB, non diventa \texttt{E:} ma una \emph{cartella}\\
  (solitamente in \texttt{/media} o \texttt{/mnt}).
\item
  \textbf{Case Sensitivity}\\
  \texttt{File.txt} ≠ \texttt{file.txt}\\
  👉 Fai fare un test pratico: è uno degli errori più comuni.
\end{itemize}

\subsection{Tour Guidato (Live Demo)}\label{tour-guidato-live-demo}

Apri il terminale e naviga con \texttt{cd}, mostrando queste directory:

\begin{itemize}
\item
  \textbf{\texttt{/bin} \& \texttt{/sbin}}\\
  I binari (eseguibili).\\
  Qui vivono comandi come \texttt{ls}, \texttt{cp}, \texttt{ip}.
\item
  \textbf{\texttt{/etc}}\strut \\
  Configurazione del sistema.\\
  Mnemonico: \textbf{Editable Text Configuration}\\
  Qui vive l'anima del sistema.
\item
  \textbf{\texttt{/home}}\strut \\
  Le cartelle degli utenti standard.
\item
  \textbf{\texttt{/root}}\strut \\
  La home dell'amministratore (spesso inaccessibile senza permessi).
\item
  \textbf{\texttt{/var}}\strut \\
  Dati variabili: log, spool di stampa, database, siti web.
\item
  \textbf{\texttt{/tmp}}\strut \\
  File temporanei (di solito si svuota al riavvio).
\item
  \textbf{\texttt{/dev}}\strut \\
  I file device (es. \texttt{/dev/null}, il buco nero).
\end{itemize}

\begin{center}\rule{0.5\linewidth}{0.5pt}\end{center}

\section{2. Editing di Testo: Nano vs Vim (30
min)}\label{editing-di-testo-nano-vs-vim-30-min}

\textbf{Obiettivo:} Configurare server senza interfaccia grafica.

\subsection{Nano (La scelta sicura)}\label{nano-la-scelta-sicura}

Editor immediato per principianti.\\
I comandi sono visibili in basso.

\subsubsection{Comandi base}\label{comandi-base}

\begin{itemize}
\tightlist
\item
  \texttt{nano\ file.txt} → apre/crea il file
\item
  \texttt{CTRL\ +\ O} → salva (Write Out)
\item
  \texttt{CTRL\ +\ X} → esce
\end{itemize}

\subsection{Vim (Il ``Mostro Finale'' --
Opzionale)}\label{vim-il-mostro-finale-opzionale}

Mostralo solo se c'è tempo o curiosità.\\
È \textbf{modale} (Insert vs Command).

\subsubsection{Panic Button}\label{panic-button}

Se qualcuno entra per sbaglio e non sa uscire:

\subsection{ESC -\textgreater{} :q! -\textgreater{}
INVIO}\label{esc---q---invio}

\section{3. Utenti e Superpoteri: sudo (45
min)}\label{utenti-e-superpoteri-sudo-45-min}

\textbf{Obiettivo:} Capire la gestione multi-utente e la sicurezza.

\subsection{Concetti}\label{concetti}

\begin{itemize}
\item
  \textbf{Root}\\
  L'amministratore supremo.
\item
  \textbf{User}\\
  L'utente limitato (quello che usano loro).
\item
  \textbf{Sudo}\\
  \emph{SuperUser DO} → permesso temporaneo per agire come root.
\end{itemize}

\subsection{Demo: I file degli utenti}\label{demo-i-file-degli-utenti}

Gli utenti sono solo righe di testo in un file.

\begin{itemize}
\item
  \texttt{cat\ /etc/passwd}\strut \\
  Lista utenti pubblica.\\
  Spiegare i campi:\\
  \textbf{User -- UID -- GID -- Home -- Shell}
\item
  \texttt{cat\ /etc/shadow}\strut \\
  Password hashate (file privato).
\end{itemize}

\subsubsection{Prova pratica}\label{prova-pratica}

\begin{Shaded}
\begin{Highlighting}[]
\FunctionTok{cat}\NormalTok{ /etc/shadow}
\end{Highlighting}
\end{Shaded}

\textbf{Permesso negato}

\begin{Shaded}
\begin{Highlighting}[]
\FunctionTok{sudo}\NormalTok{ cat /etc/shadow}
\end{Highlighting}
\end{Shaded}

→ \textbf{Funziona}

\begin{center}\rule{0.5\linewidth}{0.5pt}\end{center}

\section{4. Permessi e Proprietà (1 ora --
Cruciale)}\label{permessi-e-proprietuxe0-1-ora-cruciale}

\textbf{Obiettivo:} Decifrare \texttt{ls\ -l}.

\subsection{\texorpdfstring{La Triade:
\texttt{rwx}}{La Triade: rwx}}\label{la-triade-rwx}

Usa \texttt{ls\ -l} su un file e spiega la prima colonna\\
(es. \texttt{-rwxr-xr-\/-}).

\subsubsection{Struttura}\label{struttura}

\begin{itemize}
\item
  \textbf{Tipo}

  \begin{itemize}
  \tightlist
  \item
    \texttt{-} file
  \item
    \texttt{d} directory
  \end{itemize}
\item
  \textbf{User (u)}\\
  Proprietario del file
\item
  \textbf{Group (g)}\\
  Gruppo assegnato
\item
  \textbf{Others (o)}\\
  Tutti gli altri
\end{itemize}

\subsubsection{Significato dei permessi}\label{significato-dei-permessi}

\begin{itemize}
\tightlist
\item
  \textbf{r (Read)}

  \begin{itemize}
  \tightlist
  \item
    File: leggere\\
  \item
    Directory: listare contenuto
  \end{itemize}
\item
  \textbf{w (Write)}

  \begin{itemize}
  \tightlist
  \item
    File: modificare\\
  \item
    Directory: creare/cancellare file
  \end{itemize}
\item
  \textbf{x (Execute)}

  \begin{itemize}
  \tightlist
  \item
    File: eseguire\\
  \item
    Directory: entrarci (\texttt{cd})
  \end{itemize}
\end{itemize}

\subsection{Comandi Lab}\label{comandi-lab}

\subsubsection{\texorpdfstring{\texttt{chmod} (Change
Mode)}{chmod (Change Mode)}}\label{chmod-change-mode}

\begin{itemize}
\tightlist
\item
  \textbf{Simbolico}
\end{itemize}

\begin{Shaded}
\begin{Highlighting}[]
\FunctionTok{chmod}\NormalTok{ u+x script.sh}
\end{Highlighting}
\end{Shaded}

\begin{itemize}
\tightlist
\item
  \textbf{Ottale}
\end{itemize}

\begin{Shaded}
\begin{Highlighting}[]
\FunctionTok{chmod}\NormalTok{ 755 script.sh }\CommentTok{\# script standard}
\FunctionTok{chmod}\NormalTok{ 644 file.txt }\CommentTok{\# file standard}
\end{Highlighting}
\end{Shaded}

⚠️ \textbf{Warning:} evitare \texttt{777} se non per test disperati.

\subsubsection{\texorpdfstring{\texttt{chown} (Change
Owner)}{chown (Change Owner)}}\label{chown-change-owner}

\begin{Shaded}
\begin{Highlighting}[]
\FunctionTok{sudo}\NormalTok{ chown utente:gruppo file}
\end{Highlighting}
\end{Shaded}

\begin{center}\rule{0.5\linewidth}{0.5pt}\end{center}

\section{5. Gestione Pacchetti: APT (45
min)}\label{gestione-pacchetti-apt-45-min}

\textbf{Obiettivo:} Installare software in modo professionale.

\subsection{Concetti}\label{concetti-1}

\begin{itemize}
\item
  \textbf{Repository}\\
  Il ``magazzino'' ufficiale del software.
\item
  \textbf{Dipendenze}\\
  APT scarica automaticamente le librerie necessarie.
\end{itemize}

\subsection{Il Flusso Sacro}\label{il-flusso-sacro}

\begin{Shaded}
\begin{Highlighting}[]
\FunctionTok{sudo}\NormalTok{ apt update}
\end{Highlighting}
\end{Shaded}

Aggiorna la \textbf{lista} dei pacchetti (non il software).

\begin{Shaded}
\begin{Highlighting}[]
\FunctionTok{sudo}\NormalTok{ apt upgrade}
\end{Highlighting}
\end{Shaded}

Aggiorna i pacchetti installati.

\begin{Shaded}
\begin{Highlighting}[]
\FunctionTok{sudo}\NormalTok{ apt install nome}
\end{Highlighting}
\end{Shaded}

Installa un programma (es. \texttt{htop}, \texttt{neofetch}).

\begin{Shaded}
\begin{Highlighting}[]
\FunctionTok{sudo}\NormalTok{ apt remove nome}
\end{Highlighting}
\end{Shaded}

Rimuove il programma (mantiene i file di config).

\begin{Shaded}
\begin{Highlighting}[]
\FunctionTok{sudo}\NormalTok{ apt purge nome}
\end{Highlighting}
\end{Shaded}

Rimuove tutto.

\subsection{Lab ``Gratificazione
Immediata''}\label{lab-gratificazione-immediata}

Installare:

\begin{itemize}
\tightlist
\item
  \texttt{cmatrix} → effetto Matrix
\item
  \texttt{sl} → locomotiva se sbagliano \texttt{ls}
\item
  \texttt{htop} → monitor risorse colorato
\end{itemize}

\begin{center}\rule{0.5\linewidth}{0.5pt}\end{center}

\section{6. Note Mentali per
l'Istruttore}\label{note-mentali-per-listruttore}

\begin{itemize}
\item
  \textbf{Il TAB è tuo amico}\\
  Insistere sull'autocompletamento.\\
  Scrivere i percorsi a mano = frustrazione garantita.
\item
  \textbf{Root vs Permessi}\\
  Root ignora i permessi.\\
  Un file \texttt{chmod\ 000} può essere letto da root.
\item
  \textbf{Disclaimer \texttt{rm\ -rf}}\\
  Quando spieghi la cancellazione ricorsiva con \texttt{sudo},\\
  fai un avviso enorme:\\
  \emph{``Da grandi poteri derivano grandi responsabilità.''}
\end{itemize}

\begin{center}\rule{0.5\linewidth}{0.5pt}\end{center}

\section{7. Link e Risorse Extra (per gli
studenti)}\label{link-e-risorse-extra-per-gli-studenti}

\begin{itemize}
\item
  \textbf{Chmod Calculator}\\
  Per visualizzare e capire i permessi ottali.
\item
  \textbf{Linux FHS Official}\\
  Documentazione ufficiale della gerarchia filesystem.
\item
  \textbf{Vim Adventures}\\
  Gioco per imparare Vim divertendosi.
\end{itemize}

\end{document}
